% Possible perceptual challenges
% \todo{reflect on the results of the whole. Would I use this? pull the results together}
% \todo{talk about the difference with consonants, talk about what i would have done differently}
%Describe next steps
\section{Discussion}
\paragraph{Real-world use.} We note that \visowel{} is a vowel-specific interactive tool. It still struggles with certain consonants which can limit its real-world use. Future work may address this limitation by using additional linguistic properties in conjunction with our methods. Additionally, users could also be given greater freedom on which vowel should be targeted for practice.

Any incorrect feedback will lead users astray, which is likely to be detrimental over an extended period of time. Based on offset times, \visowel{} currently gives faulty feedback. Another possible effect of using \visowel{} for a longer time is that it could put undue emphasis on vowels. While initial interactions are encouraging since participants continued to think about pronunciation as a whole, there were many more comments regarding vowels. 

\paragraph{Longer tutorial.} 
We noticed a tendency for participants to trust the visual feedback over their own ears. Because there is no statistical way to capture just the vowel and exclude all of a consonant, introducing learners to the concept of sonorant consonants should help reduce confusion regarding longer lines on the chart. To further reduce excess visual information, we could simplify vowel lines, potentially allowing toggle between an average location and the whole vowel.

Given two participants success with anchoring vowels based on English vowels, we believe that including more information during practice explaining how to interact with \visowel{} could help. This might include more English words with different vowels to provide more reference points, and feedback to emphasize the physical connection to the chart by giving relational feedback, e.g. "Your vowel is lower than the goal, try bringing your tongue up to get closer". 

% Participants also mentioned that the spelling of some words reminded them of English sounds. We will hide the words initially so that participants can focus on hearing the sounds and not see conflicting visual information.

One of the takeaways from participant responses to \visowel{} is the need for greater transparency on how the system extracts their vowels. We received multiple comments regarding the confusion of what affected the position on the chart, one being that participants thought the pitch of their voice might make the line show up in different places. The tutorial could have practice that contrasts loud production with soft and high with low pitched speech.
% Participants were also unsure what caused the system to be unable to catch their vowels, with some assuming it might be because they were speaking too loudly. 

% Participants expressed other aspects of the practice that were difficult for them. P2 said they found the vowel chart the most confusing part of practice since they weren't aware that they could click on a line and see where the line began and ended. They suggested putting an arrow on one end to indicate the end of the vowel. \visowel{} was also confusing because it was hard to understand how the tongue moves as it produces speech. Two participants found the inclusion of resonant consonants on the chart as confusing. P3 added that they found the extra lines distracting and were unable to get their speech to different areas of the chart. Similarly, P7 said they often were not sure what to do in either practice scenario and "went off vibes".

\paragraph{Longitudinal study.}
Our goal is to improve intelligibility in a second language, not attain native-likeness. Our visualization, though an improvement over frequency markings and addressing previous limitations, does not provide visual markers sufficient for beginner learners. Based on the feedback, and to de-emphasize native-like pronunciation, incorporating an intelligibility threshold would be helpful for a second language learners. 

Another limitation to our findings is the novelty effect; none of the participants had seen the visualization before. However, since participants mentioned using the Spanish speaker's visual as a goal better explain why participants greater engagement with \visowel{}. A study over a longer period of time would allow us to understand how much of the engagement was due to novelty.

\paragraph{Applicability for other languages.}
\visowel{} can be extended to other languages, provided that they primarily differ from the first language in tongue location. We could add to the current visuals with additional language specific contexts, such as differences in lip rounding or nasalization. For future work, we could create a catalog of language specific idiosyncrasies, which could be loaded based on user needs.